% !PRINTED 446
\weberpage{446}
%<<CONT
is satisfied for
\[
u = (k), \quad v = (lk) \quad \text{and} \quad v = (l^{-1}k).
\]
But from what was proved at the end of § 69 [formula
(32)], the congruence follows from this:
\weq{16}{[(k)^p - (lk)] [(lk)^p - (k)] \equiv 0 \quad (\mathrm{mod}\ p).}
If now $\mathfrak P$ is a prime ideal that divides $p$, then one of the
two factors on the left side of (17) % [sic] German print has (17); only (16) has two factors on its left side
must be divisible by $\mathfrak P$.
Which of the two assumptions thus possible we pursue further
is immaterial, since the one passes over into the other when
$k$ is interchanged with $l^{-1}k$ and $l$ with $l^{-1}$. So
let
\weq{17}{(k)^p \equiv (lk) \quad (\mathrm{mod}\ \mathfrak P).}
If the congruence (17) is satisfied for any one of the conjugate
values $(k)$, then it holds for every other one as well.
For by § 121:
\[
(k') = f(k, \sqrt{D}), \quad (lk') = f(lk, \sqrt{D}),
\]
hence
\[
(k')^p \equiv f(k^p, \sqrt{D}) \equiv f(lk, \sqrt{D}) \equiv (lk') \quad (\mathrm{mod}\ \mathfrak P).
\]
It therefore follows from (14), % [sic] German print has (14); the substitution described applies to (17)
if $k$ is replaced by $l^{-1}k$:
\weq{18}{(l^{-1}k)^p \equiv (k) \quad (\mathrm{mod}\ \mathfrak P),}
that is, according as the one or the other of the two classes $l$
belonging to $p$ is put for $l$, the one or the other
factor of (16) is divisible by $\mathfrak P$.
By repeated application of (17) there results, for any
arbitrary positive exponent $\pi$,
\weq{19}{(k)^{p^\pi} \equiv (l^\pi k) \quad (\mathrm{mod}\ \mathfrak P)}
(inference from $\pi$ to $\pi + 1$).
If now $l$ belongs to the exponent $\varepsilon$, or $\varepsilon$ is the index of
$p$, then by (19):
\weq{20}{(k)^{p^\varepsilon} \equiv (k) \quad (\mathrm{mod}\ \mathfrak P),}
and $p^\varepsilon$ is the \gesperrt{lowest} power of $p$ that satisfies this
condition. For if a still lower power of $p$ satisfies the congruence (19), % [sic] German print cites (19) then (20); their logical roles appear interchanged
it follows from (20) that two \gesperrt{different}
class invariants $(k)$, $(k')$ are congruent with respect to the modulus $\mathfrak P$.
Their difference $(k) - (k')$ would then be divisible by $\mathfrak P$, and consequently the
discriminant of $H_m$ divisible by $p$, contrary to the hypothesis.
